% Generated draft: requires mathematical and source review.
\documentclass[11pt]{article}
\usepackage[margin=25mm]{geometry}
\usepackage{fontspec}
\setmainfont{Latin Modern Roman}
\usepackage{amsmath,amssymb,mathtools,unicode-math}
\setmathfont{Latin Modern Math}
\usepackage{enumitem,needspace,xurl,hyperref}
\setlength{\parindent}{0pt}
\setlength{\parskip}{5pt}
\setlength{\emergencystretch}{3em}
\newcommand{\N}{\mathbb{N}}
\newcommand{\Z}{\mathbb{Z}}
\newcommand{\Q}{\mathbb{Q}}
\newcommand{\R}{\mathbb{R}}
\newcommand{\C}{\mathbb{C}}
\newcommand{\F}{\mathbb{F}}
\newcommand{\A}{\mathbb{A}}
\newcommand{\PP}{\mathbb{P}}
\newcommand{\E}{\mathbb{E}}
\newcommand{\eps}{\varepsilon}
\newcommand{\dd}{\,\mathrm{d}}
\begin{document}
\section*{UnsolvedMath ID: 30001221}
\textbf{Source problem code:} catalogue entry
\subsection{Definitions and mathematical statement}
The supplied record concerns a type\(\text{-}B\) Hecke algebra and a root-of-unity specialization denoted by \(\xi\). Its rank, coefficient field, Hecke parameters, weight function, and specialization map are not specified. The following therefore records the intended assertions as an incomplete statement schema, rather than a fully specified theorem or open problem.

For a nonnegative integer \(n\), a bipartition of \(n\) is an ordered pair \(\lambda=(\lambda^{(1)},\lambda^{(2)})\) of partitions with \(|\lambda^{(1)}|+|\lambda^{(2)}|=n\). Write \(\Lambda_n\) for their set. Here \(n\) is the unspecified rank parameter of the intended algebra. Retain the source notation \(S^\xi_\lambda\) and \(D^\xi_\lambda\). Their precise constructions are absent from the record; in particular, the cellular datum needed to identify these modules must be recovered before interpreting the assertions as a theorem. For a finite-length module \(M\) and a simple module \(L\), \([M:L]\) denotes the multiplicity of \(L\) in a composition series of \(M\).

Let \(\mathcal U\subseteq\Lambda_n\) denote the intended set of Uglov bipartitions. This set depends on root-of-unity and charge data that the record does not provide. The symbol \(\mathbf s\) below retains the source superscript \(s\); its entries, admissibility conditions, and relation to the Hecke parameters are unspecified. Likewise, the record does not define the coefficients \(d^{\mathbf s}_{\lambda\mu}(v)\) or their normalization.

Subject to the missing parameter hypotheses and the source's assumptions described as Lusztig's unequal parameter conjectures \((\mathrm P1)\)--\((\mathrm P15)\), the assertions appearing in the record have the schematic form
\[
\forall\lambda\in\Lambda_n,\qquad
D^\xi_\lambda\ne0
\quad\Longleftrightarrow\quad
\lambda\in\mathcal U,
\]
and
\[
\forall\lambda\in\Lambda_n\ \forall\mu\in\mathcal U,\qquad
[S^\xi_\lambda:D^\xi_\mu]
=d^{\mathbf s}_{\lambda\mu}(1).
\]
The second display presupposes that the nonzero \(D^\xi_\mu\) are the intended simple modules. The source additionally requires labels of nonzero composition factors to be ordered by the Kazhdan--Lusztig cellular order. Neither the defining cellular datum nor the direction of the required comparison between \(\lambda\) and \(\mu\) is supplied, so no precise order implication can be stated from this record alone.

The supplied statement-verification note reports that the official report presents these assertions as a theorem under Lusztig's Conjecture 14.2 and its parameter hypotheses, rather than as an open problem. The record does not reproduce that conjecture, the parameter hypotheses, or their precise relationship to \((\mathrm P1)\)--\((\mathrm P15)\). Its statement status is \texttt{unrecoverable}. Consequently, this schema does not identify a separate unresolved mathematical assertion.
\subsection{Short English statement}
The record describes conditional claims identifying nonzero simple-module labels with Uglov bipartitions, expressing decomposition multiplicities by the coefficients \(d^{\mathbf s}_{\lambda\mu}(1)\), and restricting composition factors by a Kazhdan--Lusztig cellular order. Essential parameters and definitions are missing. Supplied verification says the report states a conditional theorem, so a precise open problem cannot be recovered from this record.
\subsection{Sources}
\begin{itemize}
\item Supplied local UnsolvedMath registry record, permanent ID \(30001221\), exact source code \texttt{catalogue entry}, entitled ``Uglov Bipartitions and Hecke Decomposition Numbers.'' The record gives the year 2009 and supplies no individual proposer attribution.
\item Supplied authoritative report citation: \emph{Representations of Finite Groups} (2009), Oberwolfach Reports, pp.\ 917--970, DOI 10.4171/owr/2009/17. Source URL and literature-source link labelled ``original report'': \url{https://doi.org/10.4171/owr/2009/17}. The supplied background describes Oberwolfach Reports as documenting mathematical workshops and their open-problem discussions. No theorem number or specific page for these assertions is supplied.
\item Supplied literature-source link labelled ``rank three results'': \url{https://www.sciencedirect.com/science/article/pii/S0021869320304865}. No author, title, or precise theorem from this linked article is supplied.
\item Literature assessment checked on \texttt{2026-08-21}: the supplied curation labels the record ``partially solved'' and states that the reported theorem is conditional on Lusztig's unequal parameter conjectures, with later work settling only important classes such as rank three. This is supplied dated curation, not independently verified here; the background itself requests independent verification of time-sensitive status and cited claims.
\item Supplied statement-verification note: the official report states the claims as a theorem under Lusztig's Conjecture 14.2 and its parameter hypotheses, rather than as an open problem. The record retains the original formulation for provenance, promotes no separate clean statement, and marks the statement \texttt{unrecoverable}. This verification note is supplied evidence and has not been independently checked here.
\end{itemize}
\end{document}
